This work studies methods for writing content into, and reading content out of, the internal workspace of frozen language models. Such methods are dual-use. On the beneficial side, the doctrine developed here is oriented toward transparency: verified, budgeted, auditable placement of nameable content, with an explicit measure of how much an intervention constrains a model's behaviour, properties useful for auditing and for studying model internals. On the risk side, activation steering could in principle be used to manipulate deployed models. Several empirical findings in this paper are directly relevant to that risk assessment: unconstrained flooding, which requires no doctrine and no insight, already dominates raw content-surfacing on small models, so the marginal misuse capability contributed by the doctrine is limited; and naive contrastive directions failed to move safety-relevant behaviour on an aligned model under registered criteria, indicating that alignment-breaking via simple steering is not demonstrated here. The recognition-gated hardening result points the other way, toward defence: it shows that the same activation signal can strengthen a model against jailbreaks with a predictable, deployable operating point on some models. That study also carries its own cautionary finding, that naive unconditional hardening can make a model less safe (attack success rising from 0.583 to 0.917 on one model), which argues for the offline clean-gap check before any such defence is deployed. The attack corpora and extraction code used to build the defence are kept in a separate, unreleased sibling library for that reason. The released steering runtime operates only on open-weight models the user already controls; the live-inference gateway is deliberately not publicly deployed. All released numbers trace to committed gate records, which the author considers a precondition for responsible claims in this area.
